% ------------------------------------------------------------
% AC 2 — IBM0515 Matemática Discreta — VERSÃO B
% Compila com pdfLaTeX
% ------------------------------------------------------------
\documentclass[11pt,addpoints,a4paper]{exam}

% ------------------------ Pacotes-base -----------------------
\usepackage[T1]{fontenc}
\usepackage[utf8]{inputenc}
\usepackage[brazil]{babel}
\usepackage{newpxtext,newpxmath}
\let\Bbbk\relax
\usepackage{microtype}
\usepackage{geometry}
\usepackage{amsmath}
\usepackage{xcolor}
\usepackage{enumitem}
\usepackage[hidelinks]{hyperref}
\usepackage{tcolorbox}
\usepackage{icomma}
\usepackage{tikz}
\usetikzlibrary{arrows.meta,positioning,calc}

% ---------------------- Geometria & espaços ------------------
\geometry{a4paper, left=2cm, right=2cm, top=2cm, bottom=1.8cm}
\setlength{\parindent}{0pt}
\setlength{\parskip}{0.4em}
\setlist[enumerate]{topsep=0.1em,itemsep=0.1em}

% ------------------------- Paleta Ibmec ----------------------
\definecolor{ibmecBlue}{HTML}{003764}
\definecolor{ibmecBlueLight}{HTML}{2A5B84}
\definecolor{ink}{HTML}{222222}
\colorlet{heading}{ibmecBlue}

% ---------------------- Cabeçalho/Rodapé ---------------------
\firstpageheader{}{}{}
\runningheader{IBM0515 - Matemática Discreta}{AC 2 - Versão B}{Prof. Dr. Ulisses A. M. Neto}
\runningheadrule
\footer{}{{\thepage} de {\numpages}}{}

% ------------------- Caixas reutilizáveis --------------------
\tcbset{boxsep=4pt,arc=3pt}
\newtcolorbox{RefBox}{colframe=heading,colback=white,boxrule=0.7pt}

% ---------------------- Questões (exam) ----------------------
\pointpoints{ponto}{pontos}
\pointformat{\textcolor{heading}{(\thepoints)}}
\renewcommand{\questionlabel}{\textbf{\color{heading}\thequestion.}}
\setlength{\fboxrule}{0.5pt}

% Caixa de resolução com altura fixa e largura dada
\newcommand{\resposta}[2][\linewidth]{%
  \fbox{\begin{minipage}[t][#2][t]{\dimexpr#1-2\fboxsep-2\fboxrule\relax}\hfill\end{minipage}}}

% ----------------------- Estilos TikZ (grafos) ---------------
\tikzset{
  v/.style={circle,fill=ink,inner sep=1.7pt},
  dir/.style={-{Stealth[length=2.2mm]},thick},
  loopA/.style={loop,out=135,in=45,min distance=11mm,looseness=14},
  loopR/.style={loop,out=45,in=-45,min distance=11mm,looseness=14},
  grafo/.style={line width=0.6pt},
  acel/.style={draw,minimum height=9.5mm,minimum width=14mm,inner sep=0pt,outer sep=0pt},
  lcel/.style={draw,minimum height=6.5mm,minimum width=7mm,inner sep=0pt,outer sep=0pt,font=\small},
  ptr/.style={-{Stealth[length=2mm]},semithick}
}
% Lista de adjacência (array de ponteiros + listas encadeadas):
% \adjlist{1/{2,3},2/{1,3},3/{}}  -> linha i / {vizinhos}; {} = ponteiro nulo
\newcommand{\adjlist}[1]{%
\begin{tikzpicture}[grafo]
\foreach \i/\vals [count=\r] in {#1} {
  \pgfmathsetmacro{\y}{-0.95*\r}
  \node[acel] (A\i) at (0,\y) {};
  \node[anchor=east] at ([xshift=-3mm]A\i.west) {\i};
  \ifx\vals\empty
    \fill (A\i.center) circle (1.3pt);
  \else
    \xdef\prev{A\i}
    \foreach \v [count=\k] in \vals {
      \pgfmathsetmacro{\x}{1.9*\k+0.6}
      \node[lcel] (N\i\k) at (\x,\y) {\v};
      \node[lcel] (P\i\k) at (\x+0.7,\y) {};
      \draw[ptr] (\prev.center) -- (N\i\k.west);
      \xdef\prev{P\i\k}
    }
    \fill (\prev.center) circle (1.3pt);
  \fi
}
\end{tikzpicture}}

% --------------------------- Documento -----------------------
\begin{document}

\begin{center}
{\LARGE IBM0515 - Matemática Discreta}\\[2pt]
{\Large AC 2 - 2026.2 \quad $\cdot$ \quad Versão B}
\end{center}

\begin{RefBox}
\textbf{Nome:} \hrulefill \quad \textbf{Matrícula:} \makebox[4cm]{\hrulefill}\\[8pt]
\textbf{Data:} \makebox[3.5cm]{\hrulefill} \hfill \textbf{Nota:} \makebox[2.5cm]{\hrulefill}
\end{RefBox}

\begin{questions}

% ===================== Questão 1 — matriz e lista (não direcionado)
\question[2]
\begin{minipage}[t]{0.46\linewidth}
Considere o grafo da figura ao lado.
\begin{enumerate}[label=(\alph*),leftmargin=*]
\item Escreva a matriz de adjacência do grafo.
\item Desenhe a lista de adjacência do grafo.
\end{enumerate}
\end{minipage}\hfill
\begin{minipage}[t]{0.4\linewidth}
\centering
\begin{tikzpicture}[grafo,baseline=(current bounding box.north)]
\node[v,label=left:1] (1) at (0,1.1) {};
\node[v,label=left:2] (2) at (1.7,2.1) {};
\node[v,label=below:3] (3) at (1.5,0) {};
\node[v,label=right:4] (4) at (3.4,1.4) {};
\node[v,label=right:5] (5) at (3.4,-0.1) {};
\draw (1) -- (2);
\draw (1) -- (3);
\draw (1) to[bend right=45] (3);
\draw (2) -- (4);
\draw (3) -- (4);
\draw (3) -- (5);
\draw (4) -- (5);
\path (2) edge[loopA] (2);
\end{tikzpicture}
\end{minipage}

\makeemptybox{\stretch{1}}

% ===================== Questão 2 — lista -> grafo direcionado e matriz
\question[2]
Considere o grafo direcionado representado pela lista de adjacência a seguir.
\begin{enumerate}[label=(\alph*)]
\item Desenhe o grafo.
\item Escreva a matriz de adjacência do grafo.
\end{enumerate}
\begin{minipage}[c]{0.4\linewidth}
\centering
\scalebox{0.8}{\adjlist{1/{3,5},2/{1,2},3/{},4/{3,3},5/{2,4}}}
\end{minipage}\hfill
\begin{minipage}[c]{0.58\linewidth}
\resposta{4.6cm}
\end{minipage}

% ===================== Questão 3 — planaridade (representação)
\question[2]
Decida se o grafo a seguir é ou não planar.

\begin{minipage}[c]{0.4\linewidth}
\centering
\begin{tikzpicture}[grafo]
\foreach \n in {1,...,8} {
  \pgfmathsetmacro{\ang}{112.5-45*(\n-1)}
  \node[v,label=\ang:\n] (\n) at (\ang:1.5) {};
}
\draw (1)--(2)--(3)--(4)--(5)--(6)--(7)--(8)--(1);
\draw (1)--(4); \draw (2)--(7); \draw (3)--(6); \draw (5)--(8);
\end{tikzpicture}
\end{minipage}\hfill
\begin{minipage}[c]{0.58\linewidth}
\resposta{4.6cm}
\end{minipage}

\newpage

% ===================== Questão 4 — Euler (desigualdades)
\question[2]
Utilizando as consequências da fórmula de Euler para grafos planares (e não o teorema de Kuratowski), prove que cada um dos grafos a seguir é não planar.

{\centering
\begin{tikzpicture}[grafo]
% ---------- (a) ----------
\foreach \n/\a in {1/90,2/30,3/-30,4/-90,5/-150,6/150}
  \node[v,label=\a:\n] (a\n) at (\a:1.35) {};
\node[v,label={[inner sep=1pt]120:7}] (a7) at (120:0.45) {};
\draw (a1)--(a2)--(a3)--(a4)--(a5)--(a6)--(a1);
\foreach \n in {1,...,6} \draw (a7)--(a\n);
\draw (a1)--(a3); \draw (a2)--(a4); \draw (a3)--(a5); \draw (a4)--(a6);
\node at (0,-2.15) {(a)};
% ---------- (b) ----------
\begin{scope}[xshift=7cm]
\foreach \n in {1,...,8} {
  \pgfmathsetmacro{\ang}{90-45*(\n-1)}
  \node[v,label=\ang:\n] (b\n) at (\ang:1.35) {};
}
\foreach \x/\y in {1/2,1/3,1/8,2/4,2/6,2/7,3/4,3/6,4/5,5/6,5/7,6/8,7/8} \draw (b\x)--(b\y);
\node at (0,-2.15) {(b)};
\end{scope}
\end{tikzpicture}\par}

\makeemptybox{\stretch{1}}

% ===================== Questão 5 — isomorfismo
\question[2]
Decida se os dois grafos a seguir são isomorfos. Se forem, dê uma função que estabelece o isomorfismo; se não forem, explique por quê.

{\centering
\begin{tikzpicture}[grafo]
% ---------- (a) ----------
\foreach \n/\a in {1/90,2/30,3/-30,4/-90,5/-150,6/150}
  \node[v,label=\a:\n] (\n) at (\a:1.4) {};
\draw (1)--(2)--(3)--(4)--(5)--(6)--(1);
\draw (1)--(4); \draw (2)--(5); \draw (3)--(6);
\node at (0,-2.3) {(a)};
% ---------- (b) ----------
\begin{scope}[xshift=5.4cm,yshift=-1.25cm]
\node[v,label=above:$b$] (b) at (0,2.5) {};
\node[v,label=above:$a$] (a) at (1.6,2.5) {};
\node[v,label=above:$e$] (e) at (3.2,2.5) {};
\node[v,label=below:$d$] (d) at (0,0) {};
\node[v,label=below:$f$] (f) at (1.1,0) {};
\node[v,label=below:$c$] (c) at (3.2,0) {};
\draw (b)--(a)--(e);
\draw (d)--(f)--(c);
\draw (b)--(d); \draw (e)--(c); \draw (a)--(f);
\draw (b)--(c); \draw (e)--(d);
\node at (1.6,-1.05) {(b)};
\end{scope}
\end{tikzpicture}\par}

\makeemptybox{\stretch{1}}

\end{questions}

\end{document}
